\section{Related Work}
\label{sec:related}

\subsection{Federated Learning and Statistical Heterogeneity}
Federated learning (FL) trains a shared model from decentralized data without moving raw examples to a central server. FedAvg established the standard pattern of local client optimization followed by server aggregation~\cite{mcmahan2017communication}, but its behavior can degrade when client data are non-IID or when clients perform unequal amounts of local work. FedProx addresses part of this instability by adding a proximal regularizer that discourages local models from drifting too far from the global iterate~\cite{li2020federated}. Momentum-based methods offer an alternative stabilization mechanism that accelerates convergence under client heterogeneity~\cite{cheng2024momentumbenefits,zaccone2025communication}. FedNova instead normalizes local updates to reduce objective inconsistency caused by heterogeneous numbers of local steps~\cite{10.5555/3495724.3496362}. Control-variate and representation-sharing methods further study client drift and personalization under heterogeneous data~\cite{pmlr-v119-karimireddy20a,collins2021exploiting}. These methods provide important optimization tools for FL, but they usually assume that each participating client can train and communicate a full model update. That assumption is restrictive for cloud--edge settings where clients may only run shallow model segments. Beyond optimization-level heterogeneity, FL deployments must also contend with statistical heterogeneity across classes~\cite{10660707,nguyen2026towards}, limited computational resources on edge devices~\cite{MAL-083,11180918}, and communication cost at scale~\cite{9084352,9760729,nguyen2022hcfl,nguyen2025federated}.

\subsection{Hierarchical and Split Federated Learning}
Hierarchical FL introduces intermediate edge servers between clients and the cloud, reducing aggregation latency and distributing the communication burden across a client-edge-cloud topology~\cite{liu2020client,9488756,rana2022hierarchical}. Split learning takes a different route by partitioning a neural network and moving the deeper layers to a server, so clients transmit intermediate activations instead of executing the whole model~\cite{vepakomma2018split,GUPTA20181}. SplitFed combines split learning with federated aggregation to improve parallelism across clients~\cite{thapa2022splitfed}, while hierarchical split-FL variants extend split execution to mobile edge systems~\cite{quan2023hiersfl}. Sub-model deployment strategies such as HeteroFL further adapt model capacity to heterogeneous client resources~\cite{diao2021heterofl}, and resource-efficient variants reduce client-side computation through kernel sampling or asymmetric private/public partitioning~\cite{niu2023overcoming,10655220}. These approaches reduce client-side computation, but they retain a coupled optimization graph: forward activations and backward gradients must move across tiers for sample-level training. This coupling is a persistent limitation: recent SFL and hierarchical-SFL work continues to study how to reduce the resulting communication and storage overhead \cite{11091500,zhang2025amperecommunicationefficienthighaccuracysplit}, mitigate the privacy leakage of transmitted smashed data \cite{qiu2023evaluating}, and preserve privacy under heterogeneity-aware masking \cite{wang2025towards}. In contrast, \hsfp decouples tier-wise learning by sending class-wise prototype statistics upward, avoiding cross-tier backpropagation and reducing dependence on synchronized activation exchange.

\subsection{Prototype-Based and Contrastive Federated Learning}
Prototype-based FL communicates class-level feature summaries rather than full parameter updates or sample activations. FedProto uses class prototypes to align heterogeneous local representations across clients~\cite{tan2022fedproto}, while class-prototype regularization has also been used to reduce non-IID effects~\cite{dai2023tacklingFedProx}. Even with class-prototype regularization, prototype drift under extreme heterogeneity remains a known failure mode~\cite{bensiah2025feddpa,zhang2020personalized}, motivating aggregation and consistency mechanisms that explicitly account for prototype reliability and staleness. Related contrastive FL methods, including MOON, constrain local representation drift by comparing local, global, and previous model representations~\cite{9578660}. Knowledge-transfer methods such as FedDF and FedGen use ensemble distillation or generator-based synthesis to transfer information at the server without direct access to private data~\cite{lin2020ensemble,zhu2021data}. Other prototype-communication schemes such as FedNTProto refine prototype-based personalization with metric-learning losses ~\cite{10660707}, and gradient-conflict-aware federated methods such as FOUL address related on-server optimization challenges in the prototype/gradient-sharing regime~\cite{nguyen2026computation}. These methods motivate communication through representations and distributions, but most are designed for flat client-server FL. They do not directly analyze a hierarchical split setting in which edge and cloud tiers train on synthetic feature distributions induced by lower-tier class-wise prototype statistics.

\subsection{Asynchronous Federated Learning and Stale-Update Convergence}
Asynchronous FL removes the requirement that the server wait for all selected clients before aggregation. Asynchronous Federated Optimization studies delayed updates in parameter space~\cite{xie2019asynchronous}, FedBuff buffers asynchronous updates before applying them~\cite{nguyen2022federated}, and general analyses characterize convergence with heterogeneous and delayed client updates~\cite{fraboni2023general}. Closest to our convergence setting, recent work directly analyzes convergence for hierarchical split federated learning under staleness and system optimization constraints~\cite{10980018}. These results are directly relevant to stale prototype updates and partial prototype updates, but the object being delayed in \ehsfp is different. Rather than aggregating stale model parameters or gradients, upper tiers train from class-wise prototype statistics whose age changes the synthetic feature distribution itself. This shifts the key stability question from delayed parameter averaging to prototype-induced distributional drift.

\subsection{Event-Driven and Serverless ML}
Serverless FL systems show how event-driven execution and elastic cloud resources can support scalable federated workloads. FedLess studies secure serverless orchestration for FL~\cite{grafberger2021fedless}, while LIFL develops a lightweight event-driven serverless platform for FL~\cite{qi2024lifl}. In this manuscript, serverless computing is not treated as the algorithmic contribution. It is the deployment motivation for stateless learning episode design: workers may be invoked episodically, receive incomplete packets, and terminate before a persistent optimizer state is available. The learning problem is therefore to stabilize upper-tier optimization under missing, stale, and partial prototype arrivals.

Overall, existing methods typically assume at least one of three properties: persistent workers, synchronized aggregation rounds, or sample-level activation/gradient exchange across split tiers. None of these assumptions fully matches prototype-decoupled hierarchical split-federated learning under stale prototype updates, partial prototype updates, missing prototype packets, and asynchronous edge/cloud aggregation. \ehsfp addresses this gap by formulating stateless learning episode training over time-varying prototype-induced distributions and by analyzing convergence under episodic prototype memory, reliability-aware prototype aggregation, prototype dropout regularization, prototype replay consistency, and residual prototype generator mechanisms.
